\section{Conclusion}
\label{sec:conclusion}

\begin{itemize}
  \item \emph{Routing.} The body's convective boundary layer takes a slow release and lets a fast one pass, so the
  momentum of a release next to a body decides which nose it reaches. On a sofa the silent release rises behind the
  releaser to the ceiling and misses both noses on its first pass; in an elevator it is strongest for the person
  standing behind the releaser. \gap{v4: standing pair}
  \item \emph{Peak bound.} Under linear transport a slower release of the same gas never gives a higher peak, in any
  room or flow. Only the near-field pre-dilution of a fast jet lets a silent release win: by a median factor of 1.45
  at equal mass in the analytic model, often below the just-noticeable ratio. As emitted, a silent passage is the
  stronger one at the neighbour in 48\% of pairs in the analytic model and at a bystander in \gap{v4; preprint: 21\%
  (0--70\%)} of pairs in the simulated scenes.
  \item \emph{Selection.} A silent passage is noticed only if it is smelled, and this selection alone can explain the
  belief: in literature observer models a remembered silent passage seems the stronger one in \gap{v4; preprint:
  68--93\%} of comparisons. The explanation predicts that releasers hold the belief less firmly than bystanders and
  that people without a sense of smell never remember a silent passage.
\end{itemize}
